\section{Motivation und Problemstellung}
\setauthor{\firstauthor}

Die im vorherigen Abschnitt beschriebene Ausgangssituation zeigt, dass der Lernprozess bei der Vorbereitung für die Matura 
stark von eigenständigem Üben und individueller Organisation geprägt ist. Gleichzeitig fehlen jedoch geeignete digitale 
Werkzeuge, die Lernenden dabei helfen, ihren Wissensstand realistisch einzuschätzen und Schwächen gezielt zu verbessern. 
Daraus ergibt sich eine zentrale Fragestellung: 
\textbf{Wie kann der Lernprozess so unterstützt werden, dass Lernen strukturierter, transparenter und nachhaltiger wird?}

Ein wesentliches Problem besteht darin, dass Fehler aktuell nur selten bewusst analysiert werden. Erhalten Lernende kein 
unmittelbares Feedback, so ist oft schwer zu klären, welche Denk- oder Verständnisfehler vorliegen und wie man diese in Zukunft
vermeiden kann. Dadurch wiederholen sie Inhalte oft unsystematisch, statt gezielt an den Gebieten zu arbeiten, in denen immer
noch Unsicherheit vorliegt. Dies erschwert den Aufbau eines langfristig stabilen Wissensfundaments.

Hinzu kommt, dass der individuelle Lernfortschritt zu wenig sichtbar ist. Für Schüler:innen bestehen kaum Möglichkeiten, ihren 
aktuellen Lernfortschritt einschätzen oder Entwicklungen über einen längeren Zeitraum nachvollziehen zu können. Besonders bei
der Maturavorbereitung führt dies zu einem Gefühl von mangelnder Kontrolle über den Lernprozess. Anstelle von strukturierter 
Vorbereitung für die Matura haben Lernende oft das Gefühl, sie würden einfach auf gut Glück lernen.

Aufgrund dieser Situation entstand die Motivation, sich im Zuge dieser Diplomarbeit mit der Frage auseinanderzusetzen, wie 
Lernprozesse im schulischen Rahmen gezielt unterstützt und verbessert werden können. Besonders wichtig ist dabei, wie digitale 
Werkzeuge einen großen Beitrag leisten können, Fehler zu identifizieren, mehr Kontrolle über den Lernprozess zu fördern und 
eigenständiges Lernen langfristig zu stärken.
